% 02-background.tex -- draft.
\section{Background: Service Mesh Data Planes}
\label{sec:background}

A service mesh has a \emph{control plane}, which distributes identities,
certificates, and routing policy, and a \emph{data plane}, which handles
application traffic. Our measurements concern the data plane, whose
architecture determines how many proxies a request crosses and where their work
is executed. \Cref{tab:configs} summarizes the configurations we study.

\paragraph{Per-pod sidecars}
In Istio's sidecar mode, an Envoy proxy is injected into every pod, and
iptables rules redirect all inbound and outbound traffic through it. A request
between two services therefore traverses two proxies, each of which terminates
TCP, performs mTLS, and parses HTTP or gRPC to apply L7 policy and collect
telemetry. Linkerd uses the same pattern with \texttt{linkerd2-proxy}, a
purpose-built proxy written in Rust that implements a smaller feature set.

\paragraph{Istio Ambient}
Ambient mode removes the sidecar. A per-node proxy, ztunnel, captures pod
traffic and carries it between nodes over mTLS in an HTTP CONNECT-based tunnel
(HBONE). ztunnel provides L4 functions only: identity, encryption, and L4
authorization. When L7 features are required, traffic is additionally routed
through a \emph{waypoint}, a shared Envoy deployment that serves a namespace or
a service account. Ambient therefore has two operating points, L4-only (IA4)
and L4 plus waypoint (IA7), which we measure separately.

\paragraph{Cilium service mesh}
Cilium implements networking, load balancing, and L3/L4 policy with eBPF
programs attached to the kernel datapath. Transparent encryption between nodes
uses WireGuard (or IPsec), and L7 policy is enforced by redirecting the matching
traffic to an Envoy instance shared by all pods on the node. Unlike the other
meshes, Cilium encrypts at the node level rather than per workload, and only
traffic subject to an L7 rule is sent through a proxy.

\paragraph{Where the cost goes}
These designs differ in three ways that matter for performance: the number of
user-space proxies on the request path (two for sidecars, zero to two for
Ambient and Cilium), whether those proxies are dedicated to a pod or shared by
a node, and how much work is moved into the kernel. Shared proxies amortize
memory but concentrate CPU on the node that hosts them, which, as we show in
\cref{sec:rq1}, matters for tail latency under load.

\begin{table}[t]
  \centering
  \caption{Configurations measured. All meshes run with mutual authentication
           and encryption enabled.}
  \label{tab:configs}
  \footnotesize
  \setlength{\tabcolsep}{4pt}
  \begin{tabular}{@{}lll@{}}
    \toprule
    ID  & Configuration            & Data plane on the path \\
    \midrule
    B   & No mesh (Cilium CNI)     & none \\
    IS  & Istio sidecar            & Envoy in every pod \\
    IA4 & Istio Ambient, L4        & per-node ztunnel \\
    IA7 & Istio Ambient + waypoint & ztunnel + per-namespace Envoy \\
    LK  & Linkerd                  & \texttt{linkerd2-proxy} in every pod \\
    CI  & Cilium service mesh      & eBPF, WireGuard, per-node Envoy \\
    \bottomrule
  \end{tabular}
\end{table}
